\documentclass[11pt]{article}
\usepackage[T1]{fontenc}
\usepackage{lmodern,amsmath,amssymb,amsthm,mathtools}
\usepackage[margin=1in]{geometry}
\usepackage{microtype,tikz-cd}
\usepackage[numbers]{natbib}
\usepackage[colorlinks=true,linkcolor=blue,citecolor=blue,urlcolor=blue]{hyperref}
\hypersetup{pdftitle={Maximality of the Dosen--Petric Boolean Category (v12)},pdfauthor={Florian Lengyel}}
\renewcommand{\theHequation}{\thesection.\arabic{equation}}
\newtheorem{theorem}{Theorem}[section]
\newtheorem{lemma}[theorem]{Lemma}
\newtheorem{proposition}[theorem]{Proposition}
\newtheorem{corollary}[theorem]{Corollary}
\newcommand{\B}{\mathsf B}
\newcommand{\Occ}{\operatorname{Occ}}
\newcommand{\Rel}{\mathbf{Rel}_{\mathrm{fin}}}
\newcommand{\FormulaSet}{\operatorname{Form}}
\title{Maximality of the Do\v{s}en--Petri\'c Boolean Category}
\author{Florian Lengyel\thanks{Email: \href{mailto:flengyel@gradcenter.cuny.edu}{\texttt{flengyel@gradcenter.cuny.edu}}}}
\date{October 31, 2026\quad(v12)}
\begin{document}
\maketitle

\begin{abstract}
We prove that the zero-identity Boolean category of Do\v{s}en and
Petri\'c freely generated by a nonempty set of propositional variables
is maximal. Adjoining an equation between distinct parallel arrows
and closing under the arrow operations and uniform formula
substitution identifies all parallel arrows. This disproves the
conjecture in Section~14.2 of the September 2007 revision of
\emph{Proof-Theoretical Coherence}. Boolean Coherence and selection of
occurrences reduce the added equation to a contextual equation on a
signed literal. Substitution gives this equation for every formula.
Two isomorphisms listed by Do\v{s}en and Petri\'c then yield equality
of the product projections. The Lean formalization proves coherence,
maximality and the model corollary for the generated presentation,
as well as the substitution equations in tagged and finite-ordinal
coordinates.
\end{abstract}

\medskip
\noindent\textit{Mathematics Subject Classification (2020).}
Primary 03G30; Secondary 03F07, 18A32.

\smallskip
\noindent\textit{Keywords.} Boolean categories, maximality, categorial proof theory,
coherence, substitution congruences.

\section{Introduction}

Do\v{s}en and Petri\'c conjecture that their freely generated
zero-identity Boolean category is not maximal
\cite[\S14.2, p.~322]{DP}. We prove that it is maximal. Any equation
between parallel arrows that does not already hold, when closed under
the arrow operations and uniform formula substitution, identifies
every pair of parallel arrows. The category before adjoining the
equation is not a preorder: for any propositional variable $p$, its
two projections $p\wedge p\to p$ are distinct.

By Boolean Coherence, distinct parallel arrows have different
relations between occurrences of propositional variables. We select
an edge belonging to one relation and not the other. Precomposition
and postcomposition isolate this edge and give the equation used in
Do\v{s}en and Petri\'c's relative-maximality argument
\cite[\S9.7, pp.~215--216]{DP}. Two isomorphisms listed in
\cite[\S12.5, p.~294]{DP} then give equality of the product projections.
The proof of \cite[Lemma~7]{DPIsomorphic} also supplies these
isomorphisms.

For cartesian categories, Do\v{s}en and Petri\'c prove maximality by
using coherence, substitution and composition to derive equality of
the product projections \cite[\S5]{DPCartesian}.
Sections~\ref{sec:isolation} and~\ref{sec:collapse} give the reduction
for the zero-identity Boolean category.

Do\v{s}en and Petri\'c also prove maximality for lattice categories
and relative maximality for dicartesian categories
\cite[\S\S8--10]{DPBicartesian}, without assuming distribution.
Their Lemma~6 in \cite{DPIsomorphic} constructs an arrow whose
relation links just two prescribed occurrences of the same variable,
provided the implication between its source and target is a tautology.
In \cite[\S16]{DPSplit}, they isolate a pair belonging to one relation
and not the other to prove maximality of the presentations of split
preorders, split equivalence relations and binary relations.

Other definitions of Boolean category start with monoidal operations
for conjunction and disjunction. Lamarche and Stra{\ss}burger
construct free graphical, $\Delta$--$\nabla$-strong
$K^\flat$- and $K^\sharp$-autonomous categories from proof nets \cite[Theorem~3.5.4]{LSFree}.
Stra{\ss}burger studies equations on monoid and comonoid structures
in $*$-autonomous categories \cite{Strassburger}.
F\"uhrmann and Pym use an order on arrows to interpret cut reduction;
their Theorems~6.4 and~6.5 give ordered soundness and completeness
\cite{FP}. The negation equations differ from those used here.
With their monoidal unit isomorphisms, the composite in
\eqref{eq:zero-identity} equals $1_A$ by their negation triangle
\cite[(11)]{FP}; in Do\v{s}en and Petri\'c's category it is $0_A$.

We use the occurrence coordinates of \cite[\S9]{Notation}: occurrences
in a compound formula form tagged disjoint unions, and substitution
replaces each occurrence by its own block of occurrences.
Section~\ref{sec:occurrences} describes these blocks, and
Appendix~\ref{app:substitution} gives the commutative diagrams relating
them to Do\v{s}en and Petri\'c's ordinal substitution operation.
Section~\ref{sec:formalization} states the Lean theorems and their
hypotheses.
Appendix~\ref{app:cut-errata} records two typographical corrections to the
cut-elimination proof in \cite{DP}.
References to \cite{DP} use the printed
pagination of its September 2007 revised version.

\section{The Boolean proof category}\label{sec:category}

Fix a nonempty set $P$ of propositional variables. Write $\FormulaSet(P)$
for the formulas built from $P$, the constants $\top,\bot$, the binary
connectives $\wedge,\vee$, and negation $\neg$. Let $\B(P)$ be the
freely generated zero-identity Boolean category defined in
\cite[\S\S14.2--14.3]{DP}. Its objects are these formulas, and its
arrows are the arrow terms of that presentation modulo its equations.
We write $\B$ when $P$ is fixed. We recall the structure and equations
used in the proof.

The connectives $\wedge$ and $\vee$ give binary products and coproducts;
$\top$ is terminal and $\bot$ is initial. For $A_1=A$ and $A_2=B$,
denote the projections and injections by
\[
 \pi_i^{A,B}:A\wedge B\longrightarrow A_i,
 \qquad
 \iota_i^{A,B}:A_i\longrightarrow A\vee B
 \qquad(i=1,2).
\]
For $f:X\to A$ and $g:X\to B$, write
$\langle f,g\rangle:X\to A\wedge B$ for their pairing.
For $f:A\to X$ and $g:B\to X$, write
$[f,g]:A\vee B\to X$ for their copairing. Write
$!_A:A\to\top$ and $\kappa_A:\bot\to A$ for the terminal and initial
arrows, and $1_A$ for the identity of $A$. Composition $g\circ f$
applies $f$ first. The symbols $f\wedge g$ and $f\vee g$ denote the
arrows induced by the product and coproduct bifunctors.

The distributor is an arrow
\[
 d_{A,B,C}:A\wedge(B\vee C)\longrightarrow(A\wedge B)\vee C.
\]
Its first-injection equation is
\begin{equation}
 d_{A,B,C}\circ(1_A\wedge\iota_1^{B,C})
   =\iota_1^{A\wedge B,C}.
 \label{eq:first-injection}
\end{equation}
This is the coproduct-injection equation of
\cite[\S11.1, p.~234]{DP}.

Negation is represented by a functor
$\neg:\B^{\mathrm{op}}\to\B$, with specified natural
double-negation and De Morgan isomorphisms. In particular, write
$n_A:\neg\neg A\to A$ for the double-negation isomorphism.
The presentation also has arrows
$\eta_A:\top\to\neg A\vee A$ and
$\epsilon_A:A\wedge\neg A\to\bot$.
The \emph{zero identity} is the endomorphism $0_A:A\to A$ defined by
\begin{equation}
 0_A=[\kappa_A,1_A]\circ(\epsilon_A\vee1_A)\circ d_{A,\neg A,A}
     \circ(1_A\wedge\eta_A)\circ\langle1_A,!_A\rangle.
 \label{eq:zero-identity}
\end{equation}
The defining zero-identity equation asserts that
\begin{equation}
 f\circ0_A=0_C\circ g
 \qquad(f,g:A\to C).
 \label{eq:zero-law}
\end{equation}
Equation~\eqref{eq:zero-identity} is $(0\eta\epsilon)$ of
\cite[\S14.2, p.~317]{DP}. Its product and coproduct unit isomorphisms
are expressed here in pairing and copairing notation.
Equation~\eqref{eq:zero-law} is the zero-identity law $(0I)$;
see \cite[\S14.3, p.~322]{DP}.

\noindent\textit{Remark.}
Appendix~\ref{app:structural-diagrams} gives diagrammatic forms of
equations~\eqref{eq:first-injection}--\eqref{eq:zero-law}, with the
source and target of every arrow displayed.

Whenever $\B(A,C)$ is nonempty, define $0_{A,C}=f\circ0_A$ for any
$f:A\to C$. Equation~\eqref{eq:zero-law} makes this independent of
$f$: both $f\circ0_A$ and $g\circ0_A$ equal $0_C\circ g$.
Taking $f=g$ gives $f\circ0_A=0_C\circ f$. Consequently zero arrows
absorb composition whenever the composites are defined:
$h\circ0_{A,C}=0_{A,D}$ for $h:C\to D$, and
$0_{A,C}\circ k=0_{X,C}$ for $k:X\to A$.

Two arrows are \emph{parallel} when they have the same source and
target. A \emph{proof congruence} on $\B$ is an equivalence relation
on each hom-set preserved by composition and every arrow operation
of the presentation, including the derived operations of pairing,
copairing, and negation.
A \emph{substitution congruence} is a proof congruence also preserved
by every uniform substitution $\sigma:P\to\FormulaSet(P)$.
Here substitution acts on the entire arrow term, including its source
and target; this is rule~$(su)$ of \cite[\S2.7, p.~51]{DP}.

The \emph{total congruence} identifies every parallel pair. A category
is \emph{thin}, or a \emph{preorder}, when each hom-set has at most one
element. We call $\B$ \emph{maximal} when its only substitution
congruences are equality and the total congruence.

\begin{theorem}\label{thm:main}
Let $P$ be a nonempty set of propositional variables. The freely
generated zero-identity Boolean category $\B(P)$ of Do\v{s}en and
Petri\'c is maximal. Equivalently, every substitution congruence on
$\B(P)$ that identifies two distinct parallel arrows identifies all
parallel arrows.
\end{theorem}

Sections~\ref{sec:isolation} and~\ref{sec:collapse} give the proof.

\section{Occurrences and substitution}\label{sec:occurrences}

For a formula $A$, write $\Occ(A)$ for the finite set of positions
occupied by propositional variables in $A$. Repeated occurrences of the
same variable occupy distinct positions. We use tagged coordinates for
these positions.

For finite sets $X,Y$, let
$X\sqcup Y=(\{1\}\times X)\cup(\{2\}\times Y)$ be their tagged
disjoint union. The occurrence sets are defined recursively by
\[
 \begin{aligned}
 \Occ(p)&=\{*\},&
 \Occ(\top)&=\Occ(\bot)=\varnothing,\\
 \Occ(\neg A)&=\Occ(A),&
 \Occ(A\mathbin\star C)&=\Occ(A)\sqcup\Occ(C)
       \quad(\star\in\{\wedge,\vee\}).
 \end{aligned}
\]
The label map $\lambda_A:\Occ(A)\to P$ assigns to each occurrence
its underlying variable. Distinct positions may have the same label.
Numbering the occurrences from left to right gives a bijection
$\Occ(A)\to\{1,\ldots,|\Occ(A)|\}$; the latter set is empty when
$|\Occ(A)|=0$. Appendix~\ref{app:ordinals} gives the substitution
formula in these finite-ordinal coordinates.

Let $\Rel$ be the category of finite sets and binary relations.
A relation $R:X\to Y$ is a subset of $X\times Y$; identities are
diagonal relations, and composition is relational composition.
Transporting the occurrence interpretation of \cite{DP} to the
tagged coordinates gives a functor
$G:\B\to\Rel$ with $GA=\Occ(A)$.
For example,
\[
 G\pi_1^{A,C}=\{((1,x),x):x\in\Occ(A)\},\qquad
 G\iota_1^{A,C}=\{(x,(1,x)):x\in\Occ(A)\}.
\]
The second projection and injection use tag~$2$. Pairing takes the
union of the two relations after tagging their targets; copairing
takes the union after tagging their sources. The relation $Gd_{A,B,C}$
is the graph of the rebracketing bijection
\[
 \Occ(A)\sqcup\bigl(\Occ(B)\sqcup\Occ(C)\bigr)
 \longrightarrow
 \bigl(\Occ(A)\sqcup\Occ(B)\bigr)\sqcup\Occ(C),
\]
which preserves each occurrence of $A,B,C$
\cite[\S2.9, p.~60]{DP}. Zero arrows have empty relations.

\begin{theorem}[Boolean Coherence, Do\v{s}en--Petri\'c]
\label{thm:coherence}
The functor $G:\B\to\Rel$ is faithful: parallel arrows $f,g$ of
$\B$ are equal if and only if $Gf=Gg$.
\end{theorem}

This is \cite[\S14.2, p.~321]{DP}. In particular, for $p\in P$, the
two projections $p\wedge p\to p$ have different occurrence
relations, so $\B$ is nonthin.

Joyal's collapse theorem concerns cartesian closed categories with a
dualizing object; see \cite[\S14.3, pp.~324--326]{DP}. The category
$\B$ is not cartesian closed. Indeed, faithfulness gives at most one
arrow $\top\to E$ for every $E$, since $\Occ(\top)$ is empty,
whereas $1_p\ne0_p$. An exponential $p\Rightarrow p$ would give
a bijection $\B(p,p)\cong\B(\top,p\Rightarrow p)$, a contradiction.

A \emph{signed literal} is $p$ or $\neg p$ for $p\in P$.
A formula is in \emph{negation normal form} when negation occurs only
immediately before a propositional variable. The double-negation and
De Morgan isomorphisms, including the isomorphisms for the truth
constants, give an isomorphism $a_A:A\to A^\circ$ to a formula in
negation normal form. The relation $Ga_A$ is the graph of a bijection
$\Occ(A)\to\Occ(A^\circ)$
\cite[\S14.1, pp.~313--315]{DP}.

For arrows between formulas in negation normal form, every edge of
the occurrence relation joins occurrences of the same signed literal.
In the presentation on negation normal forms in
\cite[\S14.2, pp.~318--321]{DP}, objects are formulas built from
signed literals, conjunction, disjunction and the truth constants.
Arrow terms are generated by identities, associativity and symmetry
maps, diagonals and codiagonals, projections and injections, mix,
distribution, unit maps, zero endomorphisms, and $\eta_p,\epsilon_p$.
The zero endomorphisms and $\eta_p,\epsilon_p$ have empty relations;
each remaining generator joins equal signed literals.
Composition, $\wedge$ and $\vee$ preserve this property.
The equivalence with $\B$ preserves the occurrence interpretation,
so the assertion follows.

Substitution replaces each occurrence of a variable by a separate
block of occurrences of its substituted formula.
Appendix~\ref{app:profiles} describes these blocks and their label maps
by commutative diagrams, following \cite[\S9]{Notation}.
Let $Q$ be a set of propositional variables and
$\sigma:P\to\FormulaSet(Q)$ a substitution. If $A[\sigma]$ denotes the
result of uniform substitution, there is a bijection
\begin{equation}
 \Occ(A[\sigma])\cong D_{A,\sigma},\qquad
 D_{A,\sigma}=\coprod_{x\in\Occ(A)}
                  \Occ\bigl(\sigma(\lambda_A(x))\bigr).
 \label{eq:occ-substitution}
\end{equation}
The set on the right consists of pairs $(x,t)$, where $x$ is an
occurrence of $A$ and $t$ is an occurrence within the formula inserted there.
Induction on $A$ proves the bijection: the variable case has one
block, the constant cases have none, binary connectives form tagged
unions, and negation preserves the occurrence set. In particular,
two occurrences of the same variable give two separately tagged
copies of its substituted formula.

Put $E_\sigma=\coprod_{p\in P}\Occ(\sigma(p))$.
Define $\operatorname{pr}_D:D_{A,\sigma}\to\Occ(A)$ and
$\operatorname{pr}_E:E_\sigma\to P$ by
$\operatorname{pr}_D(x,t)=x$ and $\operatorname{pr}_E(p,t)=p$.
The function
$\widehat\lambda_{A,\sigma}:D_{A,\sigma}\to E_\sigma$ sends
$(x,t)$ to $(\lambda_A(x),t)$. Thus
\begin{equation}
\begin{tikzcd}[column sep=large,row sep=large]
 D_{A,\sigma}\arrow[r,"\widehat\lambda_{A,\sigma}"]
   \arrow[d,"\operatorname{pr}_D"']
   &E_\sigma\arrow[d,"\operatorname{pr}_E"]\\
 \Occ(A)\arrow[r,"\lambda_A"']&P.
\end{tikzcd}
\label{eq:block-square}
\end{equation}
Both composites send $(x,t)$ to $\lambda_A(x)$.
For $(x,t),(x',t)\in D_{A,\sigma}$, their images under
$\widehat\lambda_{A,\sigma}$ agree exactly when
$\lambda_A(x)=\lambda_A(x')$, even if $x\ne x'$.

The same presentation defines $\B(Q)$ when $Q$ is empty.
Substitution on arrow terms induces a functor
\[
 F_\sigma:\B(P)\longrightarrow\B(Q),\qquad
 A\longmapsto A[\sigma],\quad f\longmapsto f[\sigma].
\]
Rule~$(su)$ makes this functor well defined on arrow terms modulo the
defining equations. It preserves identities, composition, and the
specified arrow operations. For each arrow, the blocks in
\eqref{eq:occ-substitution} describe the occurrence sets of its
substituted source and target.

\section{Isolating an occurrence}\label{sec:isolation}

Put $j_A=\iota_1^{A,\top}:A\to A\vee\top$. The contextual equation
needed in the proof is
\begin{equation}
 j_A\circ\pi_1^{A,\bot}
   =0_{A\wedge\bot,A\vee\top}.
 \tag{$K_A$}\label{eq:context}
\end{equation}
The right-hand side is also
$\iota_2^{A,\top}\circ!_\bot\circ\pi_2^{A,\bot}$: its occurrence
relation is empty, so Boolean Coherence identifies it with the zero
arrow. Thus \eqref{eq:context} is the contextual equation of
\cite[\S9.7, pp.~215--216]{DP}.

\begin{lemma}[Occurrence isolation]\label{lem:isolation}
For distinct parallel arrows $f,g:A\to C$ in $\B$, there are a
signed literal $\ell$ and arrows
\[
 i:\ell\wedge\bot\longrightarrow A,
 \qquad r:C\longrightarrow\ell\vee\top
\]
such that, after possibly interchanging $f$ and $g$,
\[
 r\circ f\circ i=j_\ell\circ\pi_1^{\ell,\bot},
 \qquad
 r\circ g\circ i=0_{\ell\wedge\bot,\ell\vee\top}.
\]
\end{lemma}

\begin{proof}
First replace $f,g$ by $a_C\circ f\circ a_A^{-1}$ and
$a_C\circ g\circ a_A^{-1}$. These remain distinct, so it suffices to
construct the arrows for the negation normal forms $A^\circ,C^\circ$.
If the arrows for these normal forms are $i^\circ,r^\circ$, the required
arrows for the original formulas are
$a_A^{-1}\circ i^\circ$ and $r^\circ\circ a_C$.
We may therefore assume that $A,C$ are in negation normal form.

Faithfulness of $G$ gives an edge belonging to exactly one of $Gf,Gg$.
Interchange $f,g$ if necessary, and choose
$(x,y)\in Gf\setminus Gg$. The occurrences $x$ and $y$ have the same
signed literal $\ell$. Set $D=\ell\wedge\bot$ and $E=\ell\vee\top$.
For every formula $V$, the arrows
\[
 b_V=\kappa_V\circ\pi_2^{\ell,\bot}:D\to V,
 \qquad
 t_V=\iota_2^{\ell,\top}\circ!_V:V\to E
\]
have empty occurrence relations.

Construct $i:D\to A$ recursively along the tagged path
to $x$. At the selected literal, use $\pi_1^{\ell,\bot}$.
At a conjunction, pair the recursively constructed arrow with $b_V$
for the unselected factor, placing it in the order specified by the
tag. At a disjunction, compose the recursively constructed arrow
with the injection of the selected summand. The displayed rules for
$G$ show inductively that $Gi$ consists of the single edge from the
sole occurrence of $D$ to $x$.

Construct $r:C\to E$ dually along the path to $y$.
At the selected literal, use $\iota_1^{\ell,\top}$.
At a conjunction, first project to the selected factor and then apply
the recursively constructed arrow. At a disjunction, copair that
arrow with $t_V$ on the unselected summand. Thus $Gr$ is the single
edge from $y$ to the sole occurrence of $E$.

Relational composition now makes $G(r\circ f\circ i)$ the one-edge
relation from $\Occ(D)$ to $\Occ(E)$, whereas
$G(r\circ g\circ i)$ is empty. Boolean Coherence gives the asserted
equations.
\end{proof}

\begin{proposition}\label{prop:all-contexts}
If a substitution congruence $\theta$ on $\B$ identifies distinct
parallel arrows, then \eqref{eq:context} holds modulo $\theta$ for
every formula $A$.
\end{proposition}

\begin{proof}
Lemma~\ref{lem:isolation} and compatibility with composition give
\eqref{eq:context} for a signed literal $\ell$.
If $\ell=p$, let $\sigma$ substitute $A$ for $p$ and fix the other variables.
This yields \eqref{eq:context} for $A$: under
\eqref{eq:occ-substitution}, the literal occurrence in each endpoint
becomes the block $\Occ(A)$, and the arrow operations are preserved
by $F_\sigma$; see also Appendix~\ref{app:shared}.

If $\ell=\neg p$, substitute $\neg A$ for $p$, obtaining the equation
for $\neg\neg A$. Transport it along $n_A:\neg\neg A\to A$.
Put
$k_A=j_A\circ\pi_1^{A,\bot}:A\wedge\bot\to A\vee\top$.
The product and coproduct equations give the commutative square
\[
\begin{tikzcd}[column sep=large,row sep=large]
 A\wedge\bot\arrow[r,"k_A"]\arrow[d,"n_A^{-1}\wedge1_\bot"']
   &A\vee\top\\
 \neg\neg A\wedge\bot\arrow[r,"k_{\neg\neg A}"']
   &\neg\neg A\vee\top\arrow[u,"n_A\vee1_\top"'].
\end{tikzcd}
\]
Precomposition and postcomposition of the zero arrow give a zero
arrow by absorption. The equation for $\neg\neg A$ therefore gives
\eqref{eq:context} for $A$.
\end{proof}

\section{Contextual collapse}\label{sec:collapse}

The next lemma identifies two isomorphisms of the types listed in
\cite[\S12.5, p.~294]{DP}.

\begin{lemma}\label{lem:structural-isos}
For every formula $A$, the arrows
\[
 \begin{aligned}
 q_A&=[j_A\circ\pi_1^{A,\bot},\iota_2^{A,\top}]
      :(A\wedge\bot)\vee\top\longrightarrow A\vee\top,\\
 e_A&=1_A\wedge j_A
      :A\wedge A\longrightarrow A\wedge(A\vee\top)
 \end{aligned}
\]
are isomorphisms in $\B$.
\end{lemma}

\begin{proof}
First take $A=p\in P$. By \cite[\S12.5, p.~294]{DP}, the source
and target of $q_p$ are isomorphic. Each has one occurrence of $p$,
so the relation of any such isomorphism is the unique bijection
between these occurrence sets. The arrow $q_p$ has the same relation.
Boolean Coherence therefore identifies $q_p$ with that isomorphism.

The same cited list gives an isomorphism
$b:p\wedge p\to p\wedge(p\vee\top)$: take $n=m=1$ in
$p^n\wedge(p^m\vee\top)\cong p^{n+m}$ and choose $p^2=p\wedge p$.
The relation $Gb$ is a bijection between two-element occurrence sets.
After precomposing $b$ with the product symmetry if necessary, it
agrees with the order-preserving bijection $Ge_p$. Coherence then
identifies $e_p$ with an isomorphism. These two identifications also
follow from the occurrence-bijection construction in the proof of
\cite[Lemma~7]{DPIsomorphic}.

For arbitrary $A$, substitute $A$ for $p$ and fix the other variables.
The functor $F_\sigma$ sends $q_p,e_p$ to $q_A,e_A$ and preserves
their inverse equations.
\end{proof}

\begin{lemma}[Contextual collapse]\label{lem:collapse}
Let $\theta$ be a proof congruence on $\B$ such that
\eqref{eq:context} holds modulo $\theta$ for every formula $A$.
Then $\B/\theta$ is a preorder.
\end{lemma}

\begin{proof}
Write $\equiv$ for equality modulo $\theta$.
The quotient retains products and coproducts: their pairing,
copairing, and defining equations pass to equivalence classes.
Fix $A$, and put
$E=A\vee\top$, $S=(A\wedge\bot)\vee\top$, and
$s=\iota_2^{A,\top}:\top\to E$.
The arrows $0_{A\wedge\bot,E}$ and $s\circ!_{A\wedge\bot}$ have
empty occurrence relations and hence are equal in $\B$. Thus
\eqref{eq:context} and the coproduct equations give
\[
 q_A\equiv[s\circ!_{A\wedge\bot},s]=s\circ!_S.
\]
The inverse equations for $q_A$ hold in the quotient, so
\[
 1_E=q_A\circ q_A^{-1}
    \equiv s\circ!_S\circ q_A^{-1}=s\circ!_E.
\]
Together with $!_E\circ s=1_\top$, this makes $E$ terminal in the
quotient.

Consequently the projection
$r=\pi_1^{A,E}:A\wedge E\to A$ is invertible, with inverse
$\langle1_A,s\circ!_A\rangle$.
Lemma~\ref{lem:structural-isos} makes $e_A$ invertible too, and the
product equations give
$r\circ e_A=\pi_1^{A,A}$.
Hence $\pi_1^{A,A}$ is invertible. Its right inverse is the diagonal
$\Delta_A=\langle1_A,1_A\rangle:A\to A\wedge A$, which must
therefore be its inverse. It follows that
\[
 \pi_2^{A,A}
   \equiv\pi_2^{A,A}\circ\Delta_A\circ\pi_1^{A,A}
   =\pi_1^{A,A}.
\]
For parallel arrows $f,g:X\to A$, pairing gives
$f=\pi_1^{A,A}\circ\langle f,g\rangle
 \equiv\pi_2^{A,A}\circ\langle f,g\rangle=g$.
Thus the quotient is a preorder.
\end{proof}

\begin{proof}[Proof of Theorem~\ref{thm:main}]
Suppose that a substitution congruence $\theta$ on $\B(P)$ identifies
two distinct parallel arrows. Proposition~\ref{prop:all-contexts}
gives \eqref{eq:context} modulo $\theta$ for every formula $A$.
Lemma~\ref{lem:collapse} then makes $\theta$ total.
\end{proof}

We also obtain maximality in the sense of
\cite[\S9.3, p.~194]{DP}. An \emph{interpretation} in a zero-identity
Boolean category $\mathcal C$ assigns an object of $\mathcal C$ to
each propositional variable. It extends to a functor
$\B(P)\to\mathcal C$ preserving the specified products, coproducts,
units, distributor, negation maps, $\eta$ and $\epsilon$.

\begin{corollary}\label{cor:models}
Let $P$ be a nonempty set of propositional variables, and let
$f,g:A\to C$ be distinct arrows of $\B(P)$.
If a zero-identity Boolean category $\mathcal C$ satisfies $f=g$
under every interpretation of the propositional variables, then
$\mathcal C$ is a preorder.
\end{corollary}

\begin{proof}
Declare two parallel arrows of $\B(P)$ equivalent when all their
interpretations in $\mathcal C$ agree. This defines a substitution
congruence: composing an interpretation with $F_\sigma$ is another
interpretation. The congruence contains $f=g$, so
Theorem~\ref{thm:main} makes it total. In particular,
$\pi_1^{p,p}=\pi_2^{p,p}$ holds under every interpretation, for any
$p\in P$. Interpret $p$ as an arbitrary object $X$ of $\mathcal C$.
Its two projections $X\wedge X\to X$ are equal. For any arrows $a,b:Y\to X$, composing
$\langle a,b\rangle:Y\to X\wedge X$ with the equal projections gives $a=b$.
\end{proof}

\section{Lean formalization}\label{sec:formalization}

In the Lean~4 development \cite{Formalization}, formulas and arrow
terms are inductive types, with arrow terms indexed by their source
and target. The relation \texttt{Derives} is generated by the
presentation's equations and congruence rules. Its quotient is the
category $\B(P)$. Substitution, normalization and the occurrence
interpretation are defined on terms and proved to respect
\texttt{Derives}.

The theorem \texttt{relative\_maximality} states that a
substitution congruence identifying distinct parallel arrows is total.
It assumes the arrow operations, a faithful occurrence interpretation,
isomorphisms to negation normal forms, preservation of signed literals
along occurrence-relation edges, uniform substitution and double
negation, with their equations.
Its proof of contextual collapse uses the map
\[
 h_A=[1_A,1_A]\circ m_{A,A}:A\wedge A\longrightarrow A,
\]
where the mix arrow $m_{A,C}:A\wedge C\to A\vee C$ is given, for any
formula $B$, by
\[
 m_{A,C}=(\pi_1^{A,B}\vee1_C)\circ d_{A,B,C}
              \circ(1_A\wedge\iota_2^{B,C}).
\]
This is equation $(dm)$ of \cite[\S11.1, p.~234]{DP}.
The right-hand side is independent of $B$. The symmetry of mix gives
$h_A\circ\langle\pi_2^{A,A},\pi_1^{A,A}\rangle=h_A$.
Together with \eqref{eq:context} for all formulas, this identity yields
equality of the product projections.

Cut elimination and Boolean Coherence are proved for the presentation
defining $\B(P)$. Normalization, preservation of signed literals and
uniform substitution on terms and their source and target formulas
satisfy the remaining hypotheses of \texttt{relative\_maximality}.
Its application gives \texttt{Presentation.maximality}.

The theorem \texttt{Presentation.short\_maximality} formalizes the proof
in Section~\ref{sec:collapse}, with explicit inverses for $q_A$ and $e_A$.
The model corollary is proved using interpretations preserving the
specified operations and their composition with substitution. The formalization
also verifies the occurrence-profile and finite-ordinal constructions of
Appendix~\ref{app:substitution}, the diagram identities of
Appendix~\ref{app:structural-diagrams}, and the cut reductions and rank
comparisons discussed in Appendix~\ref{app:cut-errata}.

Five declarations are registered in the Palomar Registry
\cite{palomar-2026-10-30-000014-v1}:
the two maximality proofs, the congruence dichotomy, nonzero variable
identities, and the absence of arrows from $\top$ to $\bot$.

\appendix
\section{Substitution in occurrence coordinates}\label{app:substitution}

Do\v{s}en and Petri\'c's generalized composition of definable
connectives \cite[\S2.5, pp.~45--46]{DP} substitutes formula
patterns and computes the profile of the resulting pattern.
We describe this profile by the tagged coordinates of
\cite[\S\S9.2--9.3]{Notation}. Each occurrence after substitution
is indexed by an occurrence of the outer pattern and an occurrence
of the pattern inserted there.

\subsection{Profiles and their fibres}\label{app:profiles}

For an integer $n\geq0$, put $\underline n=\{1,\ldots,n\}$, with
$\underline0=\varnothing$. A \emph{formula pattern} on a finite set
of placeholders is a formula in those placeholders. Its
\emph{occurrence profile} records the placeholder at each position,
reading positions from left to right.

Let $U=\{p_1,\ldots,p_k\}$ consist of distinct placeholders and let
$M\in\FormulaSet(U)$ have $m$ occurrences. Its profile is
\[
 o_M:\underline m\longrightarrow\underline k,
 \qquad o_M(r)=i\quad\text{when position $r$ contains $p_i$}.
\]
For each $i\in\underline k$, choose a set
$Q_i=\{q_{i,1},\ldots,q_{i,\ell_i}\}$ of distinct local placeholders and a
pattern $N_i\in\FormulaSet(Q_i)$ with $n_i$ occurrences. The sets
$Q_i$ are pairwise disjoint. Write
$o_i:\underline{n_i}\to\underline{\ell_i}$ for the profile of $N_i$.
Empty placeholder sets and empty occurrence sets are allowed; a
pattern over an empty placeholder set has no variable occurrences.

Define the sets of occurrence coordinates and placeholder indices by
\[
 \begin{aligned}
 D&=\{(r,t):r\in\underline m,\ t\in\underline{n_{o_M(r)}}\},\\
 C&=\{(i,s):i\in\underline k,\ s\in\underline{\ell_i}\}.
 \end{aligned}
\]
Thus $(r,t)$ specifies an occurrence in the copy of $N_{o_M(r)}$
inserted at position $r$ of $M$. The pair $(i,s)$ specifies a local
placeholder of $N_i$. Define the injections
\[
 a_r:\underline{n_{o_M(r)}}\longrightarrow D,\quad a_r(t)=(r,t),
 \qquad
 b_i:\underline{\ell_i}\longrightarrow C,\quad b_i(s)=(i,s).
\]
There is a unique function $\gamma:D\to C$ for which, for every
$r\in\underline m$, the following square commutes:
\begin{equation}\label{eq:app-fibre}
\begin{tikzcd}[column sep=large,row sep=large]
 \underline{n_{o_M(r)}}\arrow[r,"a_r"]\arrow[d,"o_{o_M(r)}"']
   &D\arrow[d,"\gamma"]\\
 \underline{\ell_{o_M(r)}}\arrow[r,"b_{o_M(r)}"']&C.
\end{tikzcd}
\end{equation}
Every element of $D$ is $a_r(t)$ for a unique pair $(r,t)$.
Thus these squares determine $\gamma$ uniquely by
\[
 \gamma(r,t)=\bigl(o_M(r),o_{o_M(r)}(t)\bigr).
\]

Define the block projections $p_D:D\to\underline m$ and
$p_C:C\to\underline k$ by $p_D(r,t)=r$ and $p_C(i,s)=i$.
Then $p_C\circ\gamma=o_M\circ p_D$, so the following square commutes:
\begin{equation}\label{eq:app-profile}
\begin{tikzcd}[column sep=large,row sep=large]
 D\arrow[r,"\gamma"]\arrow[d,"p_D"']&C\arrow[d,"p_C"]\\
 \underline m\arrow[r,"o_M"']&\underline k.
\end{tikzcd}
\end{equation}
Both composites send $(r,t)$ to $o_M(r)$.
If $r\ne r'$ but $o_M(r)=o_M(r')=i$, the copies of
$\underline{n_i}$ indexed by $r$ and $r'$ are disjoint subsets
of $D$.

To identify this construction with pattern substitution,
replace each $q_{i,s}$ in $N_i$ by $(i,s)$ and denote the resulting
pattern over $C$ by $\widetilde N_i$. Let
$\Phi_0:U\to\FormulaSet(C)$ be given by $\Phi_0(p_i)=\widetilde N_i$.
Every occurrence of $M[\Phi_0]$ has unique coordinates $(r,t)\in D$.
This gives a bijection $\chi:\Occ(M[\Phi_0])\to D$, and its relation
to the label map is the commutative square
\begin{equation}\label{eq:app-occurrence-labels}
\begin{tikzcd}[column sep=large,row sep=large]
 \Occ(M[\Phi_0])\arrow[r,"\chi","\cong"']
    \arrow[d,"\lambda_{M[\Phi_0]}"']&D\arrow[d,"\gamma"]\\
 C\arrow[r,"1_C"']&C.
\end{tikzcd}
\end{equation}
Indeed, the occurrence with coordinates $(r,t)$ contains precisely
the placeholder $\bigl(o_M(r),o_{o_M(r)}(t)\bigr)$.

\subsection{The finite-ordinal presentation}\label{app:ordinals}

The finite ordinals used in \cite{DP} number the elements of each
block consecutively. Put
\[
 s_0=t_0=0,\qquad
 s_r=\sum_{q=1}^{r}n_{o_M(q)}\quad(1\leq r\leq m),\qquad
 t_i=\sum_{j=1}^{i}\ell_j\quad(1\leq i\leq k).
\]
The block encodings are the bijections
\[
 \begin{aligned}
 \theta_D:D&\longrightarrow\underline{s_m},
       &\theta_D(r,t)&=s_{r-1}+t,\\
 \theta_C:C&\longrightarrow\underline{t_k},
       &\theta_C(i,s)&=t_{i-1}+s.
 \end{aligned}
\]
Define $\delta:\underline{s_m}\to\underline{t_k}$ as the unique
function making the following square commute:
\begin{equation}\label{eq:app-ordinal}
\begin{tikzcd}[column sep=large,row sep=large]
 D\arrow[r,"\gamma"]\arrow[d,"\theta_D","\cong"']
    &C\arrow[d,"\theta_C","\cong"']\\
 \underline{s_m}\arrow[r,"\delta"']&\underline{t_k}.
\end{tikzcd}
\end{equation}
Equivalently, $\delta=\theta_C\circ\gamma\circ\theta_D^{-1}$.
For each $(r,t)\in D$, this gives
\begin{equation}\label{eq:app-coordinate}
 \delta(s_{r-1}+t)=t_{o_M(r)-1}+o_{o_M(r)}(t)
 \qquad\bigl((r,t)\in D\bigr).
\end{equation}
Every source position belongs to exactly one nonempty block, so this
formula determines $\delta$ even when some blocks are empty.

In Do\v{s}en and Petri\'c's notation, $\mu$ and $\nu_i$ correspond to
$o_M$ and $o_i$, and $\mu(\nu_1,\ldots,\nu_k)$ corresponds to $\delta$.
Their block locator selects the unique $r$ for which
$s_{r-1}<x\leq s_r$; substituting $t=x-s_{r-1}$ in
\eqref{eq:app-coordinate} gives their formula
\cite[\S2.5, p.~46]{DP}. The square therefore describes the
occurrence-profile component of their generalized composition.
The \emph{skeleton} of a pattern is obtained by replacing every
placeholder by one fixed symbol and leaving the constants and
connectives unchanged. A skeleton and an occurrence profile recover
the pattern by assigning the indicated placeholder to each position.
The generalized composite therefore consists of the skeleton of the
substituted pattern and the profile $\delta$.

For example, take
\[
 M(p,q)=q\wedge(p\vee q),\qquad
 N_1(a,b)=a\wedge b,\qquad
 N_2(c,d)=d\wedge(c\vee d).
\]
The profiles, listed by successive values, are
$o_M=(2,1,2)$, $o_1=(1,2)$, and $o_2=(2,1,2)$.
The source blocks have sizes $3,2,3$, while the target blocks list
$(a,b)$ and $(c,d)$. Thus
\[
\begin{array}{c|ccc|cc|ccc}
 \text{source position}&1&2&3&4&5&6&7&8\\\hline
 \text{placeholder}&d&c&d&a&b&d&c&d\\
 \delta&4&3&4&1&2&4&3&4
\end{array}
\]
For instance, source position $2$ has coordinates $(1,2)$.
The profile map gives $\gamma(1,2)=(2,1)$, and the target encoding
then gives $\theta_C(2,1)=3$. Thus $\delta(2)=3$.

\subsection{Shared variables and the maximality argument}\label{app:shared}

The local placeholders may be assigned names in a common variable
set $Q$. Let $\rho:C\to Q$ be that assignment; different elements
of $C$ may have the same image. Replace each placeholder $c\in C$
by the variable $\rho(c)$, obtaining the renaming map
$(-)[\rho]:\FormulaSet(C)\to\FormulaSet(Q)$. Define
$\sigma:U\to\FormulaSet(Q)$ by
$\sigma(p_i)=\widetilde N_i[\rho]$.
Then the following square commutes:
\begin{equation}\label{eq:app-renaming}
\begin{tikzcd}[column sep=large,row sep=large]
 \FormulaSet(U)\arrow[r,"{(-)[\Phi_0]}"]
    \arrow[d,"{(-)[\sigma]}"']
   &\FormulaSet(C)\arrow[d,"{(-)[\rho]}"]\\
 \FormulaSet(Q)\arrow[r,"1"']&\FormulaSet(Q).
\end{tikzcd}
\end{equation}
At $p_i$, both routes give $\sigma(p_i)$. The constant and connective
clauses agree, so structural induction on the formula proves commutativity.
Renaming changes the labels on the occurrences. Thus the same variable
of $Q$ may occur in several inserted patterns while the pairs $(r,t)$
distinguish its occurrences.

In Proposition~\ref{prop:all-contexts}, the source and target of
the contextual equation are $p\wedge\bot$ and $p\vee\top$.
Each contains one occurrence of $p$. Substituting $A$ for $p$
therefore gives one block $\Occ(A)$ at each endpoint, and the
substituted arrows have type $A\wedge\bot\to A\vee\top$. If the isolated literal is $\neg p$,
the substitution $p\mapsto\neg A$ instead gives
$\neg\neg A\wedge\bot\to\neg\neg A\vee\top$; the
double-negation square in Proposition~\ref{prop:all-contexts}
then gives the required equation for $A$.

\section{Diagrams for three defining equations}
\label{app:structural-diagrams}

\paragraph{The first-injection equation.}
For formulas $A,B,C$, equation~\eqref{eq:first-injection} is the
commutative triangle
\[
\begin{tikzcd}[column sep=large,row sep=large]
 A\wedge B
   \arrow[r,"1_A\wedge\iota_1^{B,C}"]
   \arrow[dr,"{\iota_1^{A\wedge B,C}}"']
   &A\wedge(B\vee C)\arrow[d,"d_{A,B,C}"]\\
   &(A\wedge B)\vee C.
\end{tikzcd}
\]
Both paths have source $A\wedge B$ and target $(A\wedge B)\vee C$.
The upper arrow is induced by $1_A:A\to A$ and
$\iota_1^{B,C}:B\to B\vee C$.

\paragraph{The zero identity.}
For a formula $A$, equation~\eqref{eq:zero-identity} is expressed by
the commutative diagram
\[
\begin{tikzcd}[column sep=large,row sep=large]
 A\arrow[r,"{\langle1_A,!_A\rangle}"]\arrow[d,"0_A"']
   &A\wedge\top\arrow[r,"1_A\wedge\eta_A"]
   &A\wedge(\neg A\vee A)\arrow[d,"d_{A,\neg A,A}"]\\
 A
   &\bot\vee A\arrow[l,"{[\kappa_A,1_A]}"']
   &(A\wedge\neg A)\vee A\arrow[l,"\epsilon_A\vee1_A"'].
\end{tikzcd}
\]
The path across the top, down the right side, and back along the
bottom is the five-arrow composite in~\eqref{eq:zero-identity}.
Its first arrow pairs $1_A:A\to A$ with $!_A:A\to\top$;
its last arrow copairs $\kappa_A:\bot\to A$ with $1_A:A\to A$.
The arrows between them use
$\eta_A:\top\to\neg A\vee A$ and
$\epsilon_A:A\wedge\neg A\to\bot$, together with the distributor.

\paragraph{The zero law.}
For any parallel arrows $f,g:A\to C$,
equation~\eqref{eq:zero-law} is the commutative square
\[
\begin{tikzcd}[column sep=large,row sep=large]
 A\arrow[r,"0_A"]\arrow[d,"g"']&A\arrow[d,"f"]\\
 C\arrow[r,"0_C"']&C.
\end{tikzcd}
\]
The upper-right path is $f\circ0_A:A\to C$, and the lower-left
path is $0_C\circ g:A\to C$. The arrows $f$ and $g$ need not be equal.

\clearpage
\section{Two corrections to the cut-elimination proof}\label{app:cut-errata}

The cut-elimination proof in \cite[\S11.2]{DP} contains two
typographical errors in its rank comparisons. We use the notation of
that proof: $r$ is the rank of the cut being reduced and $r'$ is the
rank of a recursive subcut. The left and right ranks are positive;
their sum is the cut rank \cite[p.~249]{DP}.

\paragraph{Page 260, case $(w\,\gamma\,3)$.}
In the sentence immediately following the displayed reduction,
the inequality $r<r'$ should read $r'<r$.
The recursive cuts have the same cut formula and smaller rank:
the corresponding one-sided ranks decrease on passing to the
indicated premises, while the other one-sided rank is unchanged.

\paragraph{Page 262, case $(w\,\vee\,1)$.}
In the continuation of this case, the phrase ``whose rank is~1''
should read ``whose right rank is~1'', as in the analogous
argument on p.~261. The selected formula in the right premise
is newly introduced by the displayed $\vee L$ rule, so its right
rank is~1. Its total rank is the unchanged left rank plus~1,
and is therefore at least~2.

\section*{Use of AI tools}
OpenAI's ChatGPT assisted with the development of the proof, preparation
of the manuscript, and checks of the arguments and references. OpenAI Codex
was used to develop the Lean formalization and to run its verification
checks. The author takes responsibility for the mathematical claims and
the final text.

\clearpage
\begingroup
\small
\setlength{\bibsep}{2pt}
\bibliographystyle{unsrtnat}
\bibliography{references_v12,palomar_v12}
\endgroup
\end{document}

